% Source: https://artofproblemsolving.com/community/c3h1949218_geometry_formula_proving
% Stable identity: aops-topic-1949218-q01
\textbf{$K=rs$.} Let $I$ be the incentre. The incircle touches all three sides, so the distance from $I$ to each side is $r$.

\begin{asy}
usepackage("amsmath");
size(5.5cm);
defaultpen(fontsize(10pt));
pair A=(0,0),B=(6,0),C=(1.6,3.8);
real a=abs(B-C),bb=abs(C-A),c=abs(A-B);
pair I=(a*A+bb*B+c*C)/(a+bb+c);
pair ft(pair P,pair X,pair Y){return X+dot(P-X,Y-X)/dot(Y-X,Y-X)*(Y-X);}
real r=abs(I-ft(I,A,B));
draw(A--B--C--cycle,linewidth(1));draw(circle(I,r));
draw(I--A,gray(0.5));draw(I--B,gray(0.5));draw(I--C,gray(0.5));
draw(I--ft(I,A,B),dashed);draw(I--ft(I,B,C),dashed);draw(I--ft(I,C,A),dashed);
dot(I);label("$I$",I,1.2*dir(20));
label("$r$",(I+ft(I,A,B))/2,E);
label("$A$",A,SW);label("$B$",B,SE);label("$C$",C,N);
label("$a$",(B+C)/2,NE);label("$b$",(A+C)/2,NW);label("$c$",(A+B)/2,S);
label("\mbox{}",(min(currentpicture,user=true)+max(currentpicture,user=true))/2,black); // keeps the text after the figure black when LaTeX typesets the labels inline
\end{asy}

Joining $I$ to the vertices cuts the triangle into triangles $IBC$, $ICA$, $IAB$ with bases $a,b,c$ and common height $r$. Hence
\[
K=[IBC]+[ICA]+[IAB]=\dfrac{ar}{2}+\dfrac{br}{2}+\dfrac{cr}{2}=r\cdot\dfrac{a+b+c}{2}=rs .
\]

\textbf{$K=\dfrac{abc}{4R}$.} First we show $c=2R\sin C$. Let $O$ be the circumcentre and $M$ the midpoint of $AB$; since $OA=OB=R$, triangle $OAB$ is isosceles and $OM\perp AB$.

\begin{asy}
usepackage("amsmath");
size(5.5cm);
defaultpen(fontsize(10pt));
pair A=(0,0),B=(6,0),C=(1.6,3.8);
real d=2*(A.x*(B.y-C.y)+B.x*(C.y-A.y)+C.x*(A.y-B.y));
real ux=((A.x^2+A.y^2)*(B.y-C.y)+(B.x^2+B.y^2)*(C.y-A.y)+(C.x^2+C.y^2)*(A.y-B.y))/d;
real uy=((A.x^2+A.y^2)*(C.x-B.x)+(B.x^2+B.y^2)*(A.x-C.x)+(C.x^2+C.y^2)*(B.x-A.x))/d;
pair O=(ux,uy);real R=abs(A-O);pair M=(A+B)/2;
draw(circle(O,R));draw(A--B--C--cycle,linewidth(1));
draw(O--A,blue);draw(O--B,blue);draw(O--M,dashed);
draw(M+0.25*dir(90)--M+0.25*dir(90)+(0.25,0)--M+(0.25,0));
dot(O);label("$O$",O,N);label("$M$",M,S);
label("$R$",(O+B)/2,NE);
label("$A$",A,SW);label("$B$",B,SE);label("$C$",C,N);
label("\mbox{}",(min(currentpicture,user=true)+max(currentpicture,user=true))/2,black); // keeps the text after the figure black when LaTeX typesets the labels inline
\end{asy}

If $C$ is acute, the inscribed angle theorem gives $\angle AOB=2C$, so $\angle AOM=C$ and $c=2AM=2R\sin C$. If $C$ is obtuse, $C$ lies on the minor arc $AB$, the central angle on the other side is $\angle AOB=360^\circ-2C$, so $\angle AOM=180^\circ-C$ and again $c=2R\sin(180^\circ-C)=2R\sin C$. If $C=90^\circ$, $AB$ is a diameter and $c=2R=2R\sin C$.

Now $K=\dfrac{1}{2}ab\sin C$, and $\sin C=\dfrac {c}{2R}$, so
\[
K=\dfrac{1}{2}ab\cdot\dfrac{c}{2R}=\dfrac{abc}{4R}. \qquad\blacksquare
\]
